\documentclass[main.tex]{subfiles}
\begin{document}

\section{Exercise 6 – CMRR Improvement in Differential Amplifier}

\begin{figure}[H]
    \centering
    \includegraphics[width=0.85\linewidth]{images/exercise6_circuit.png}
    \caption{CMOS differential amplifier circuit}
    \label{fig:ex6_circuit}
\end{figure}

\subsection*{CMRR Fundamentals}

Common-Mode Rejection Ratio quantifies differential amplifier performance:
$$\text{CMRR} = \frac{|A_d|}{|A_{cm}|} \quad \text{or} \quad \text{CMRR}_{(dB)} = 20\log\left(\frac{A_d}{A_{cm}}\right)$$

For single-ended output, $A_{cm} \approx -R_D/(2R_B)$ where $R_B$ is tail source output resistance. Thus $\text{CMRR} \approx g_m R_B$ - directly proportional to tail impedance.

\textbf{Current limitations:} MB (W/L = 20/100) has long channel for moderate $r_o$, but simple current source limits CMRR. Device mismatch in M1-M2 and load M3-M4 further degrades rejection.

\subsection*{Part a) CMRR Improvement Methods}

\textbf{1. Cascode Tail Current Source:}

Replace MB with cascode structure (two stacked transistors). Output impedance increases to $r_{out} \approx g_m r_{ds}^2$, dramatically improving CMRR (20-40 dB gain typical).

\textbf{2. Improved Matching:}

Increase W/L of M1-M2 (currently 20/5) since mismatch $\sigma(\Delta V_T) \propto 1/\sqrt{WL}$. Larger devices reduce random variations. Use common-centroid layout and matched orientation.

\textbf{3. Cascode Active Load:}

Replace simple M3-M4 mirror with cascode configuration. Increases differential gain $A_d$ and output impedance, thus improving CMRR.

\textbf{4. Bias Voltage Filtering:}

Add bypass capacitor at tail bias node to AC ground high-frequency common-mode signals. Improves AC CMRR but not DC.

\subsection*{Part b) Trade-offs Analysis}

\begin{table}[H]
\centering
\small
\begin{tabular}{|p{2.5cm}|p{2cm}|p{1.8cm}|p{1.5cm}|p{2cm}|}
\hline
\textbf{Method} & \textbf{Headroom} & \textbf{Power} & \textbf{Area} & \textbf{Speed} \\
\hline
Cascode Tail & $-2V_{DS,sat}$ penalty & Same current & +1 device & Minimal impact \\
\hline
Large W/L & No penalty & Same for $I_D$ & Significant $\uparrow$ & Reduced BW (parasitics) \\
\hline
Cascode Load & Output swing $\downarrow$ & Same & +2 devices & Dominant pole $\downarrow$ \\
\hline
Bypass Cap & No penalty & No change & Large cap area & AC only \\
\hline
\end{tabular}
\end{table}

\textbf{Cascode Tail:}
\textit{Pros:} Major CMRR boost (20-40 dB), simple implementation
\textit{Cons:} Requires $2V_{DS,sat}$ headroom, reduces input common-mode range, process sensitive

\textbf{Improved Matching:}
\textit{Pros:} Better statistical matching, reduced offset, lower noise
\textit{Cons:} Significant area increase, higher parasitics reduce bandwidth

\textbf{Cascode Load:}
\textit{Pros:} Higher $A_d$ boosts CMRR, better PSRR
\textit{Cons:} Reduced output swing, needs additional bias voltages, affects stability (compensation required)

\textbf{Bypass Capacitor:}
\textit{Pros:} Simple, low cost, improves AC PSRR
\textit{Cons:} DC CMRR unchanged, large capacitor for low frequencies

\subsection*{Recommended Design Approach}

Combine cascode tail (primary CMRR improvement) with moderate W/L increase (e.g., M1-M2: 40/5 instead of 20/5). This balances CMRR, area, and headroom. Add bypass capacitor for AC noise rejection. Reserve cascode load for applications requiring maximum gain where voltage swing can be sacrificed.

\end{document}
